% Bewertungspaket zum Gesamttest «Abbildungen und Geraden im Koordinatensystem» — Quelle des PDFs (python3 scripts/build-lp-pdf.py koordinaten).
% Ganze Punkte je Teilschritt; Werte und Fehlerwerte mit python3 nachgerechnet (scripts/lp/koordinaten/zahlen.py).
% Gerüst aus dem Bewertungspaket von LP 3 (gleichungen-formeln): Auftrag an die KI, Folgefehler, (E)-Zeilen.
\documentclass[11pt]{article}
\usepackage{../lp-druck}
\usepackage{ko-gitter}
\renewcommand{\lpname}{Abbildungen und Geraden im Koordinatensystem}
\renewcommand{\lpteilgebiet}{}
\renewcommand{\lpfuss}{Bewertungspaket}
\newcolumntype{P}{>{\centering\arraybackslash}p{10mm}}
\newenvironment{raster}{\par\smallskip\noindent\renewcommand{\arraystretch}{1.6}\linespread{1.35}\selectfont\tabularx{\linewidth}{@{}>{\raggedright\arraybackslash}p{38mm}P X@{}}\toprule
  \small\textsc{Teilschritt} & \small\textsc{P} & \small\textsc{Musterlösung}\\\midrule}{\bottomrule\endtabularx\par}
\newcommand{\fehler}[1]{\par\smallskip{\small\textcolor{lprot}{\bfseries Typische Fehler:} #1}}
\newcommand{\E}{\,{\footnotesize\bfseries\color{lpgruen}(E)}}
\begin{document}
\lpkopf{Bewertungspaket zum Gesamttest}{}

\begin{lpachtung}{Erst lösen, dann öffnen}
Dieses Paket enthält die Musterlösung. Wer es vor dem Lösen liest, misst am Ende nur, wie gut das Abschreiben gelingt.
\end{lpachtung}

\section*{1 · So gehst du vor}
\begin{enumerate}[leftmargin=*,itemsep=2pt]
\item \textbf{Gesamttest lösen} — vollständig, mit Rechenweg, ohne dieses Paket und ohne Taschenrechner, mit Lineal und Geodreieck.
\item \textbf{Selbst bewerten} mit Abschnitt~3: Zeile für Zeile vergleichen und die Punkte eintragen. Das geht ohne Hilfsmittel
  und ist der Normalfall.
\item \textbf{Oder bewerten lassen.} Ob du dafür eine KI nutzt und welche, entscheidest du selbst und in eigener Verantwortung —
  je nachdem, was dir aufgrund deines Alters und deiner persönlichen Situation erlaubt ist (Altersgrenzen und
  Nutzungsbedingungen des Anbieters, allenfalls Einverständnis deiner Eltern). Fotografiere oder scanne deine Lösung gut lesbar,
  \textbf{ohne Namen}. Handyfotos tragen oft unsichtbar Standort, Zeit und Gerät mit (Metadaten) — lade lieber einen Scan oder ein
  Bildschirmfoto des Fotos hoch. Lade nur deine Lösung, das Testblatt und dieses PDF hoch, keine persönlichen Daten, und füge den Auftrag aus
  Abschnitt~2 ein.
\item \textbf{Rückmeldung prüfen.} Stimmt, was die KI aus deiner Handschrift gelesen hat? Eine KI kann sich irren —
  im Zweifel gilt das Raster in Abschnitt~3.
\item \textbf{Gezielt wiederholen} nach Abschnitt~4.
\end{enumerate}

\section*{2 · Auftrag an die KI}
\begin{tcolorbox}[colback=lplinie!25,colframe=lplinie,boxrule=0.5pt,arc=1mm,fontupper=\small]
Du bist Korrektorin oder Korrektor für Mathematik auf der Sekundarstufe I in der Schweiz. Im Anhang findest du (1)~das Testblatt mit den Aufgaben, (2)~das Bewertungspaket zum Gesamttest «Abbildungen und Geraden im Koordinatensystem» mit Musterlösung und Punkteraster und (3)~Bilder meiner handschriftlichen Lösung.\par\medskip
Geh so vor:\par
1. Schreib zu jeder Aufgabe G1 bis G8 zuerst ab, was du in meiner Lösung liest — Rechenweg, Koordinaten, Ergebnis; bei Zeichnungen (G1, G4, G5, G6), welche Punkte wo eingezeichnet sind. Was du nicht sicher lesen kannst, markierst du mit [unsicher]. Nicht raten. Achte besonders auf Vorzeichen, die Reihenfolge der Koordinaten $(x \mid y)$ und Dezimalpunkte.\par
2. Vergib die Punkte genau nach dem Raster im Bewertungspaket (Abschnitt 3), ganze Punkte je Zeile. Rechnung und Zeichnung werden getrennt bewertet: Zeichne ich meine eigenen (falschen) Bildpunkte richtig ein, gibt es den Zeichnungspunkt. Ziehe einen Fehler nur einmal ab: Rechne ich mit einem falschen Zwischenergebnis richtig weiter, gibt es die Folgepunkte. Ausnahme: Zeilen mit \textbf{(E)} gibt es nur für das richtige Ergebnis, nie als Folgepunkt. Die Zeilen «Typische Fehler» sagen, welche Rasterzeile bei diesem Fehler 0 Punkte gibt — sie sind kein zusätzlicher Abzug.\par
3. Andere richtige Wege zählen voll, zum Beispiel das Bild durch Abzählen im Gitter statt mit $2Z - P$, die Fläche in G6 mit einer anderen Zerlegung, $b$ in G7 mit dem Punkt $B$ statt $A$. Gleichwertig: $y = 3x - 2$ und $y = 3 \cdot x + (-2)$; Brüche und Dezimalzahlen mit gleichem Wert.\par
4. Begründe jeden Abzug in einem Satz und nenne die Fehlerart (zum Beispiel «$P + Z$ statt $2Z - P$», «falsche Koordinate geändert», «Zentrum vergessen», «Fläche mal k statt k²», «$\Delta x : \Delta y$ statt $\Delta y : \Delta x$», «Vorzeichenfehler»). Schreib die Musterlösung nicht ab — zeig mir, wo mein Fehler liegt.\par
5. Gib zum Schluss eine Tabelle «Aufgabe | Punkte | Begründung», die Gesamtpunktzahl (von 25) und — nach der Selbsteinschätzung in Abschnitt 4 — welches Kapitel des Leitprogramms ich wiederholen soll. Keine Note.\par
6. Fehlt ein Bild oder ist eine Aufgabe unlesbar, sag es, statt sie zu bewerten.
\end{tcolorbox}

\needspace{16\baselineskip}
\section*{3 · Musterlösung und Punkteraster}
\textbf{Allgemein:} Ganze Punkte je Zeile. Ein Fehler wird nur einmal abgezogen (Folgefehler). Gleichwertige Wege zählen voll.
In der Spalte P: \textbf{(E)}~= Ergebnispunkt, Ergebnis richtig \emph{und} Weg erkennbar, direkt aus den Angaben; ohne Zeichen = die
Zeile bewertet den Weg, die Skizze oder die Begründung. Der ganze Test ist ohne Taschenrechner gestellt. \textbf{Zeichnungszeilen} bewerten das Einzeichnen der eigenen Werte.
\textbf{Koordinaten} vertauscht ($(1 \mid -2)$ statt $(-2 \mid 1)$) gelten als falsch; fehlt nur die Einheit, gibt es keinen Abzug. \textbf{Ansatzzeilen} bewerten die richtige Gleichung; ein Rechenfehler danach kostet nur die Ergebniszeile.
«Typische Fehler» nennt, welche Zeile 0 Punkte gibt, und die Punkte, die bleiben.
\textbf{Teilrichtige Zeilen:} Verlangt eine Zeile mehrere Angaben und stimmt nur ein Teil, gibt die Zeile 0 Punkte,
ausser die Zeile nennt ausdrücklich eine Aufteilung.

\begin{lpaufgabe}{G1}{Dreieck am Punkt spiegeln}{3 P}
\begin{raster}
(a) Weg & 1 & $P' = 2Z - P$ komponentenweise, mit $2Z = (2 \mid 0)$, oder im Gitter: von $P$ bis $Z$ und gleich weit weiter\\
(a) Ergebnis & 1\E & $A'(0 \mid -1)$, $B'(-3 \mid -2)$, $C'(-1 \mid -4)$ (alle drei richtig)\\
(b) Zeichnung & 1 & die eigenen drei Bildpunkte richtig eingezeichnet und verbunden\\
\end{raster}
\fehler{$P + Z$ gerechnet ($(3 \mid 1)$, $(6 \mid 2)$, $(4 \mid 4)$): Zeilen (a) 0, (b) als Folgepunkt $\to$ 1 von 3\,P. $Z - P$ ($(-1 \mid -1)$, $(-4 \mid -2)$, $(-2 \mid -4)$): Zeilen (a) 0, (b) als Folgepunkt $\to$ 1 von 3\,P.}
\end{lpaufgabe}

\begin{lpaufgabe}{G2}{Zentrum gesucht}{3 P}
\begin{raster}
(a) Ansatz & 1 & $Z$ ist der Mittelpunkt von $P$ und $P'$: $Z = \left(\tfrac{-3 + 5}{2} \;\middle|\; \tfrac{4 + 0}{2}\right)$\\
(a) Ergebnis & 1\E & $Z(1 \mid 2)$\\
(b) & 1 & $Q' = 2Z - Q = (2 - 4 \mid 4 - 5) = (-2 \mid -1)$; als Folgepunkt mit dem eigenen $Z$\\
\end{raster}
\fehler{Halbe Differenz $\tfrac{P' - P}{2} = (4 \mid -2)$ statt Mittelpunkt: Zeilen (a) 0, (b) als Folgepunkt $\to$ 1 von 3\,P. (b) $Q + Z = (5 \mid 7)$: Zeile (b) 0 $\to$ 2 von 3\,P.}
\end{lpaufgabe}

\begin{lpaufgabe}{G3}{Drei Achsen}{3 P}
\begin{raster}
(a) $x$-Achse & 1\E & $(4 \mid 1)$: nur $y$ wechselt das Vorzeichen\\
(b) $y$-Achse & 1\E & $(-4 \mid -1)$: nur $x$ wechselt das Vorzeichen\\
(c) $y = x$ & 1\E & $(-1 \mid 4)$: die Koordinaten tauschen\\
\end{raster}
\fehler{(a) und (b) vertauscht ($(-4 \mid -1)$ an der $x$-Achse): Zeilen (a) und (b) 0 $\to$ 1 von 3\,P. (c) $(1 \mid -4)$ (getauscht und Vorzeichen gewechselt, das ist die Spiegelung an $y = -x$): Zeile (c) 0.}
\end{lpaufgabe}

\begin{lpaufgabe}{G4}{Achse gesucht}{3 P}
\begin{raster}
(a) Achse & 1 & $x = 2$, denn die Achse geht senkrecht durch den Mittelpunkt von $A$ und $A'$: $\tfrac{-2 + 6}{2} = 2$ (Gleichung und Begründung)\\
(b) & 1 & $Q'(3 \mid -1)$: $Q$ liegt $1$ links der Achse, $Q'$ $1$ rechts; $y$ bleibt. Als Folgepunkt mit der eigenen Achse\\
(c) Zeichnung & 1 & senkrechte Achse durch $x = 2$, $Q$ und das eigene $Q'$ im gleichen Abstand auf derselben Höhe\\
\end{raster}
\fehler{Achse $x = 4$ (halbe Differenz $\tfrac{6 - (-2)}{2}$): Zeile (a) 0, (b) als Folgepunkt $Q'(7 \mid -1)$, (c) als Folgepunkt $\to$ 2 von 3\,P. (b) $y$ geändert ($(1 \mid 1)$ oder $(3 \mid 1)$): Zeile (b) 0.}
\end{lpaufgabe}

\begin{lpaufgabe}{G5}{Dreieck strecken}{4 P}
\begin{raster}
(a) Weg & 1 & Weg vom Zentrum mal $k$: $P' = Z + 2 \cdot (P - Z)$, zum Beispiel $A' = (1 + 2 \cdot 1 \mid 1 + 2 \cdot 0)$\\
(a) Ergebnis & 1\E & $A'(3 \mid 1)$, $B'(7 \mid 1)$, $C'(3 \mid 5)$ (alle drei richtig)\\
(b) Zeichnung & 1 & die eigenen drei Bildpunkte richtig eingezeichnet und verbunden\\
(c) Flächen & 1 & $\tfrac{2 \cdot 2}{2} = 2$ und $\tfrac{4 \cdot 4}{2} = 8$~cm$^2$ (oder Einheitsquadrate); Faktor $4 = k^2$. Als Folgepunkt mit den eigenen Bildpunkten, wenn das Bild ähnlich ist\\
\end{raster}
\fehler{Zentrum vergessen, vom Ursprung aus gestreckt ($(4 \mid 2)$, $(8 \mid 2)$, $(4 \mid 6)$): Zeilen (a) 0, (b) und (c) als Folgepunkte $\to$ 2 von 4\,P. Faktor $2$ statt $4$ für die Fläche: Zeile (c) 0.}
\end{lpaufgabe}

\begin{lpaufgabe}{G6}{Fläche aus Koordinaten}{3 P}
\begin{raster}
Rechteck & 1 & umschliessendes Rechteck $8 \cdot 6 = 48$ (in der Zeichnung erkennbar)\\
Dreiecke & 1 & drei rechtwinklige Randdreiecke $\tfrac{8 \cdot 2}{2} = 8$, $\tfrac{5 \cdot 4}{2} = 10$, $\tfrac{3 \cdot 6}{2} = 9$\\
Ergebnis & 1\E & $48 - 8 - 10 - 9 = 21$~cm$^2$\\
\end{raster}
\fehler{Randdreiecke nicht halbiert ($48 - 16 - 20 - 18 = -6$): Zeilen Dreiecke und Ergebnis 0 $\to$ 1 von 3\,P (eine negative Fläche hätte auffallen müssen). Nur das Rechteck ($48$): Zeilen Dreiecke und Ergebnis 0 $\to$ 1 von 3\,P.}
\end{lpaufgabe}

\begin{lpaufgabe}{G7}{Geradengleichung}{3 P}
\begin{raster}
Steigung & 1 & $m = \dfrac{\Delta y}{\Delta x} = \dfrac{7 - (-5)}{3 - (-1)} = \dfrac{12}{4} = 3$\\
$b$ & 1 & einen Punkt einsetzen: $-5 = 3 \cdot (-1) + b$, also $b = -2$ (oder mit $B$: $7 = 9 + b$); Folgepunkt mit dem eigenen $m$\\
Gleichung & 1\E & $y = 3 \cdot x - 2$\\
\end{raster}
\fehler{$\tfrac{\Delta x}{\Delta y} = \tfrac{1}{3}$: Zeilen Steigung und Gleichung 0, $b$ als Folgepunkt $\to$ 1 von 3\,P. Vorzeichen bei den negativen Koordinaten verloren ($\tfrac{7 - 5}{3 - 1} = 1$): Zeilen Steigung und Gleichung 0 $\to$ 1 von 3\,P.}
\end{lpaufgabe}

\begin{lpaufgabe}{G8}{Punktprobe}{3 P}
\begin{raster}
(a) & 1 & $3 \cdot 50 - 2 = 148$: Die Gleichung stimmt, $P$ liegt auf der Geraden (Rechnung und Schluss)\\
(b) & 1 & $3 \cdot (-20) - 2 = -62 \neq -61$: $Q$ liegt nicht auf der Geraden (Rechnung und Schluss)\\
(c) & 1\E & $(-20 \mid -62)$\\
\end{raster}
\fehler{Schluss ohne Rechnung («sieht so aus», abgelesen): Zeile 0. (b) $-60 + 2 = -58$ (Vorzeichen): Zeile (b) 0, (c) $(-20 \mid -58)$: Zeile (c) 0 $\to$ 1 von 3\,P.}
\end{lpaufgabe}

\needspace{12\baselineskip}
\section*{4 · Selbsteinschätzung}
\begin{tabularx}{\linewidth}{@{}l X@{}}
\toprule
\textbf{22 – 25 P} & Die geprüften Teile sitzen. Wo du Punkte verloren hast: das Kapitel dieser Aufgabe nochmals.\\
\textbf{17 – 21 P} & Den schwächsten Teil nochmals, gemessen am Anteil (Teil A 6, B 6, C 7, D 6 P): Tüfteln und Übungen des Kapitels dieses Teils.\\
\textbf{11 – 16 P} & Zurück zu den Kapiteln aller Aufgaben, in denen du Punkte verloren hast.\\
\textbf{0 – 10 P} & Zurück zu den Kapiteln jedes Teils, in dem du weniger als die Hälfte erreicht hast; sind es alle vier, zum Vortest und zu Kapitel 1.\\
\bottomrule
\end{tabularx}

\smallskip
\textbf{Aufgabe $\to$ Kapitel:} G1, G2 $\to$ 1; G3, G4 $\to$ 2; G5, G6 $\to$ 3; G7, G8 $\to$ 4.
\end{document}
